% !TEX root = ../main.tex

\chapter{جزئیات محیط آزمایش و تنظیمات اجرا}

\section{مشخصات سخت‌افزار}

اجرای تعاملی و آزمایش‌های محلی نهایی روی سیستم زیر انجام شدند.
خروجی برخی ارزیابی‌های مقایسه‌ای مدل‌های \lr{ASR} از دفترچه‌های
منجمد \lr{Colab} به دست آمده است؛ ازاین‌رو زمان‌های هر جدول فقط
درون همان محیط و همان آزمایش مستقیماً قابل مقایسه‌اند.

\begin{table}[H]
\centering
\caption{مشخصات سخت‌افزار محیط آزمایش}
\label{tab:hardware_spec}

\begin{tabular}{ll}
\toprule
\textbf{مؤلفه} & \textbf{مشخصات} \\
\midrule
\lr{CPU} & \lr{Intel Core i7-13700H} \\
\lr{GPU} & \lr{NVIDIA GeForce RTX 4070 Laptop GPU} \\
\lr{GPU Memory} & \lr{8 GB} \\
\lr{RAM available to WSL2} & \lr{8 GB} \\
\lr{Operating System} & \lr{Linux under WSL2} \\
\lr{GPU Driver} & \lr{610.62} \\
\bottomrule
\end{tabular}
\end{table}

استفاده از یک سخت‌افزار ثابت برای مدل‌های هر گروه باعث می‌شود
مقایسه زمان اجرا و مصرف منابع تا حد امکان تحت شرایط یکسان
انجام شود.


% =========================================================
\section{محیط نرم‌افزاری}

پیاده‌سازی پروژه با زبان \lr{Python} انجام شده است.
کتابخانه‌های مورد استفاده بسته به مدل شامل ابزارهای پردازش
گفتار، مدل‌های زبانی و یادگیری عمیق هستند.

\begin{latin}
\begin{verbatim}
Python: 3.12.13
PyTorch: 2.12.0
Transformers: 4.57.6
Accelerate: 1.14.0
BitsAndBytes: 0.49.2
Vosk: 0.3.45
Piper-TTS: 1.4.2
Hazm: 0.12.1
Librosa: 0.11.0
\end{verbatim}
\end{latin}

نسخه‌های فوق محیط اجرایی جاری را مستند می‌کنند. دفترچه‌های
قدیمی‌تر ممکن است نسخه کتابخانه متفاوتی داشته باشند؛ فایل خروجی،
هش و متادیتای آزمایش‌های انتشار خطا برای جلوگیری از اختلاط
نسخه‌ها در شاخه \lr{benchmark} نگهداری شده‌اند.


% =========================================================
\section{تنظیمات ارزیابی \lr{ASR}}

مجموعه اصلی ارزیابی \lr{ASR} شامل:

\begin{center}
\lr{10,519}
نمونه صوتی
\end{center}

با مجموع مدت زمان تقریبی:

\begin{center}
\lr{52,047.28 seconds}
\end{center}

است.

برای تمام مدل‌ها معیارهای زیر ثبت شدند:

\begin{itemize}
    \item \lr{WER}
    \item \lr{CER}
    \item \lr{RTF}
    \item زمان استنتاج
    \item متوسط تأخیر
    \item مصرف حافظه \lr{GPU}
\end{itemize}

ضریب زمان واقعی از رابطه زیر محاسبه شد:

\begin{equation}
RTF =
\frac{T_{\mathrm{inference}}}
     {T_{\mathrm{audio}}}
\end{equation}


% =========================================================
\section{تنظیمات ارزیابی مدل زبانی}

ارزیابی اصلی مدل‌های زبانی شامل ۱۰۵۰ پرسش چندگزینه‌ای
\lr{ParsiNLU} و ۶۰۰ نمونه درک مطلب \lr{ParsBench RC} بود.
آزمایش ۱۰۰۰ نمونه تصادفی \lr{FarsiInstruct} فقط مطالعه مقدماتی
و خارج از رتبه‌بندی اصلی است.

معیارهای اصلی عبارت بودند از:

\begin{itemize}
    \item دقت چندگزینه‌ای؛
    \item تطابق دقیق؛
    \item \lr{F1} توکنی؛
    \item شباهت معنایی؛
    \item زمان تأخیر؛
    \item تعداد توکن در ثانیه؛
    \item و بیشینه مصرف حافظه \lr{GPU}.
\end{itemize}

برای مقایسه منصفانه، تنظیمات تولید پاسخ باید برای تمام مدل‌ها
تا حد امکان یکسان باقی بمانند.


% =========================================================
\section{تنظیمات مسیر اجرایی و ارزیابی مقاومت}

تنظیمات پیش‌فرض مسیر \lr{main.py} عبارت‌اند از:

\begin{latin}
\begin{verbatim}
ASR model: vosk-model-fa-0.42
Audio preprocessing: low_level_gain
ASR alternatives: 3
Word-confidence threshold: 0.65
Alternative-score gap: 1.0
Clarification options: 2
History turns: 6
Correction-memory items: 12
LLM: Gemma 2 9B Instruct, 4-bit
Maximum LLM input tokens: 2048
Maximum LLM output tokens: 150
TTS: Piper fa_IR-gyro-medium
\end{verbatim}
\end{latin}

آزمایش بهبود صوت و انتشار پایین‌دستی روی ۷۶۹ ضبط و ۲۸۶ شناسه
معنایی انجام شد. تقسیم \lr{N-best} شامل ۱۹۲ شناسه توسعه و ۹۴
شناسه نگه‌داشته‌شده بود. فاصله‌های اطمینان با ۵۰۰۰ تکرار
بوت‌استرپ خوشه‌ای و بذر ثابت ۲۰۲۶۰۸۲۸ محاسبه شدند. واحد
بازنمونه‌گیری شناسه معنایی بود، نه فایل ضبط منفرد.


% =========================================================
\section{تنظیمات ارزیابی \lr{TTS}}

آزمایش \lr{TTS} روی
\lr{1052}
نمونه متنی انجام شد.

برای ارزیابی کیفیت ادراکی از \lr{NISQA} و برای بررسی
قابل‌فهم‌بودن خروجی از بازشناسی مجدد صوت و محاسبه \lr{WER} و \lr{CER}
استفاده شد.

معیارهای عملکرد اجرایی نیز شامل:

\begin{itemize}
    \item زمان تولید؛
    \item متوسط تأخیر؛
    \item \lr{RTF}؛
    \item و ضریب افزایش سرعت
\end{itemize}

بودند.

ضریب افزایش سرعت از رابطه زیر به دست می‌آید:

\begin{equation}
Speedup = \frac{1}{RTF}
\end{equation}
